% Summary of curriculum experiment: best configuration per metric (mean over step lengths 25, 100, 500)
\begin{table}[H]
\centering
\small
\begin{tabulary}{\textwidth}{lLLL}
\toprule
Metric & Type & Mixture NCA (best \(NB\)) & Stochastic Mixture NCA (best \(NB\)) \\
\midrule
KL Divergence & Distributional & \(NB=5\) (plateau after \(NB=3\)) & \(NB=5\) (lowest; large gain vs.\ det.) \\
Chi-Square & Distributional & \(NB=3\) (degrades at \(NB=5\)) & \(NB=5\) (lowest) \\
Categorical MMD & Distributional & \(NB=3\) (slight rise at \(NB=5\)) & \(NB=5\) (lowest) \\
\midrule
Tumor Size Diff & Spatial & \(NB=2\) (peaks at \(NB=3\)) & \(NB=5\) (recovers after \(NB=3\) peak) \\
Border Size Diff & Spatial & \(NB=3\) (best) & \(NB=5\) (poor at \(NB=3\), then recovers) \\
Spatial Variance Diff & Spatial & \(NB=3\) (best; worsens at \(NB=5\)) & \(NB=5\) (peak at \(NB=3\), then best) \\
\bottomrule
\end{tabulary}
\caption{Curriculum experiment: summary of which neighborhood size yields the best mean value (over step lengths 25, 100, 500) for each metric and model. ``Type'' distinguishes distributional (global statistics) vs.\ spatial (geometry/structure). The deterministic model often peaks at \(NB=3\) on spatial metrics; the stochastic model peaks in error at \(NB=3\) on spatial metrics then improves at \(NB=5\).}
\label{tab:impl:curriculum:metrics_summary}
\end{table}
